% =============================================================================
% Semana 07 — Variáveis Categóricas + Revisão para Prova I
% Aula de 3 horas (19h–22h) | 25/03/2026
% =============================================================================

\chapter{Semana 07 (25/03/2026) — Variáveis Categóricas e Revisão para Prova I}

% ---------------------------------------------------------------------------
\section{Objetivo da semana}
% ---------------------------------------------------------------------------

Pessoal, esta é a última aula da Unidade I. Fechamos com a análise de
\textbf{variáveis categóricas}: tabelas de contingência, proporções
condicionais e visualizações de frequência. Depois fazemos uma
\textbf{revisão completa e integrada} de tudo que vimos desde a semana 01
para vocês chegarem confiantes na Prova I.

\begin{BoardBox}
\textbf{Roteiro da aula (19h–22h)}
\begin{enumerate}
  \item 19h–19h30 — \textbf{TEORIA:} Tabela de frequência simples para categóricas
  \item 19h30–20h — \textbf{TEORIA:} Tabela de contingência e proporções condicionais
  \item 20h–20h20 — \textbf{TEORIA:} Gráficos de barras — simples, agrupadas, empilhadas
  \item 20h20–20h50 — \textbf{PRÁTICA EM DUPLA:} Tabelas 1 e 2
  \item 20h50–21h50 — \textbf{REVISÃO:} Unidade I — resumo integrativo + simulado
  \item 21h50–22h — \textbf{Dúvidas finais e orientação para a Prova I}
\end{enumerate}
\end{BoardBox}

\begin{NoteBox}
A Prova I está prevista para \textbf{26 de março (quinta-feira)}.
Revisem a biblioteca de funções Python da turma — algumas questões pedem
que o aluno interprete saídas de código.
\end{NoteBox}

% =============================================================================
\section{Parte I — Teoria: Variáveis Categóricas (19h–20h20)}
% =============================================================================

% ---------------------------------------------------------------------------
\subsection{19h–19h30: Tabela de Frequência para Categóricas}
% ---------------------------------------------------------------------------

\begin{BoardBox}
\textbf{Frequência simples de variável categórica (para o quadro)}

Para variáveis categóricas (nominal ou ordinal), a tabela de frequência
conta \emph{quantas vezes cada categoria aparece} e calcula a proporção.

\begin{center}
\begin{tabular}{|l|r|r|}
\hline
\textbf{Categoria} & \textbf{Freq. Absoluta ($f_i$)} & \textbf{Freq. Relativa ($f_i/n$)} \\
\hline
Básico    & 120 & 40\% \\
Standard  & 105 & 35\% \\
Premium   &  75 & 25\% \\
\hline
\textbf{Total} & 300 & 100\% \\
\hline
\end{tabular}
\end{center}

\textbf{O que observar:}
\begin{itemize}
  \item Distribuição equilibrada ou há dominância de uma categoria?
  \item A proporção da menor categoria é suficiente para análise confiável?
  (categorias com $< 5\%$ são instáveis em comparações)
  \item Ordem das categorias: nominal (sem ordem natural) ou ordinal (ex.:
  ``ruim/regular/bom/ótimo'' — deve respeitar a ordem).
\end{itemize}
\end{BoardBox}

% ---------------------------------------------------------------------------
\subsection{19h30–20h: Tabela de Contingência}
% ---------------------------------------------------------------------------

\begin{SolvedBox}
\textbf{O que é uma tabela de contingência (crosstab)?}

Ao cruzar \textbf{duas variáveis categóricas}, obtemos uma tabela de
dupla entrada que mostra a distribuição conjunta. Cada célula mostra
a frequência de observações com aquela combinação de categorias.

\textbf{Exemplo:} Canal de Venda × Plano do Cliente

\begin{center}
\begin{tabular}{|l|r|r|r|r|}
\hline
& \textbf{Básico} & \textbf{Standard} & \textbf{Premium} & \textbf{Total} \\
\hline
Online & 80 & 60 & 20 & 160 \\
Presencial & 40 & 45 & 55 & 140 \\
\hline
\textbf{Total} & 120 & 105 & 75 & 300 \\
\hline
\end{tabular}
\end{center}

\textbf{Proporções condicionais por linha (leitura de decisão):}
\begin{center}
\begin{tabular}{|l|r|r|r|}
\hline
& \textbf{Básico} & \textbf{Standard} & \textbf{Premium} \\
\hline
Online & 50\% & 37,5\% & 12,5\% \\
Presencial & 28,6\% & 32,1\% & 39,3\% \\
\hline
\end{tabular}
\end{center}

\textbf{Leitura:} clientes online preferem fortemente o plano Básico (50\%).
Clientes presenciais têm distribuição mais equilibrada, com forte presença
Premium (39\%). Isso sugere que o canal influencia — ou é influenciado por —
o tipo de contratação.

\textbf{Proporções por coluna} respondem a outra pergunta: ``Dentro dos que
contratam Premium, qual canal é mais frequente?'' Depende da pergunta de negócio.
\end{SolvedBox}

% ---------------------------------------------------------------------------
\subsection{20h–20h20: Gráficos de Barras}
% ---------------------------------------------------------------------------

\begin{BoardBox}
\textbf{Quando usar barras simples, agrupadas ou empilhadas (para o quadro)}

\begin{tabular}{|l|p{5cm}|p{5.5cm}|}
\hline
\textbf{Tipo} & \textbf{Quando usar} & \textbf{Limitação} \\
\hline
Barras simples & 1 variável categórica × frequência & Não compara subgrupos \\
Barras agrupadas & 2 categorias, comparação lado a lado & Fica poluído com $>4$ grupos \\
Barras empilhadas & Composição interna de um total & Difícil comparar categorias intermediárias \\
Barras empilhadas (100\%) & Proporção relativa por grupo & Não mostra volume absoluto \\
\hline
\end{tabular}

\textbf{Regra de ouro:} barras agrupadas para \emph{comparar} categorias;
barras empilhadas 100\% para \emph{mostrar composição}.
\textbf{Pizza nunca} quando há mais de 3 fatias ou diferenças sutis — o olho humano
não consegue comparar ângulos com precisão acima de 3 fatias.
\end{BoardBox}

% =============================================================================
\section{Parte II — Prática em Dupla (20h20–20h50)}
% =============================================================================

\begin{BoardBox}
\textbf{Tabela 1 — Atendimentos por Canal e Resultado}

\medskip
\begin{tabular}{|c|l|l|l|}
\hline
\textbf{ID} & \textbf{Canal} & \textbf{Resultado} & \textbf{Turno} \\
\hline
01 & Chat    & Resolvido   & Manhã  \\
02 & Telefone & Escalado   & Tarde  \\
03 & Chat    & Resolvido   & Manhã  \\
04 & Email   & Resolvido   & Manhã  \\
05 & Telefone & Resolvido  & Tarde  \\
06 & Chat    & Escalado    & Tarde  \\
07 & Email   & Escalado    & Manhã  \\
08 & Chat    & Resolvido   & Tarde  \\
09 & Telefone & Resolvido  & Tarde  \\
10 & Email   & Resolvido   & Manhã  \\
\hline
\end{tabular}

\medskip
\textbf{Tarefas:}
\begin{enumerate}
  \item Construa a tabela de frequência de Canal (absoluta + relativa).
  \item Construa a tabela de contingência Canal × Resultado.
  \item Calcule a proporção de ``Escalado'' por canal (prop. condicional por linha).
  \item Qual canal tem maior taxa de escalamento? O que isso sugere?
\end{enumerate}

\textbf{Respostas esperadas:}
\begin{itemize}
  \item Canal: Chat 4×(40\%), Telefone 3×(30\%), Email 3×(30\%)
  \item Contingência: Chat (3R, 1E); Telefone (2R, 1E); Email (2R, 1E)
  \item Taxa de escalamento: Chat = 25\%; Telefone = 33\%; Email = 33\%
  \item Telefone e Email têm taxa de escalamento 1,3× maior que Chat.
  Sugestão: investigar se tickets escalados via Telefone/Email têm complexidade
  maior ou se é falta de capacitação no atendimento assíncrono.
\end{itemize}
\end{BoardBox}

% =============================================================================
\section{Parte III — Revisão Integrada: Unidade I (20h50–21h50)}
% =============================================================================

\begin{SolvedBox}
\textbf{Mapa mental da Unidade I — para o quadro}

\textbf{Semana 01–02:} contexto AED; ciclo de análise; tipos de variável;
tabela de frequência; importação de dados em Python.\\
\textbf{Semana 03:} medidas de posição — média, mediana, moda;
leitura crítica; quando cada uma é mais representativa.\\
\textbf{Semana 04:} medidas de dispersão — amplitude, variância, DP, IQR;
detecção de outliers (regra de Tukey); boxplot.\\
\textbf{Semana 05:} histograma e distribuições; assimetria positiva e negativa;
dados ausentes (MCAR, MAR, MNAR) e estratégias de tratamento.\\
\textbf{Semana 06:} correlação de Pearson; scatter plot; correlação $\neq$ causalidade;
variáveis confundidoras.\\
\textbf{Semana 07:} tabelas de contingência; proporções condicionais;
gráficos de barras agrupadas e empilhadas.
\end{SolvedBox}

\begin{BoardBox}
\textbf{Simulado — 10 questões (20 min para responder, 10 min para correção)}

\textbf{Q1.} Um dataset de salários tem Média = R\$12.000 e Mediana = R\$4.500.
O que isso indica sobre a distribuição? Qual medida de posição você reportaria?

\textbf{Q2.} Duas turmas têm Média = 7,0. Turma A: DP = 0,3. Turma B: DP = 3,1.
O que isso significa? Em qual turma o processo de ensino parece mais consistente?

\textbf{Q3.} Dado o conjunto $\{5, 7, 8, 10, 45\}$:
(a) Calcule $Q_1$, $Q_3$ e IQR.
(b) Determine se 45 é outlier pela regra de Tukey.

\textbf{Q4.} Um histograma de tempo de espera mostra dois picos — um em torno de
5 minutos e outro em torno de 35 minutos. O que isso sugere sobre a população analisada?

\textbf{Q5.} Você tem 15\% de missing na coluna ``renda mensal'' de uma pesquisa
de satisfação. Ao investigar, percebe que clientes de baixa renda tendem a não
responder. Que tipo de missing é esse? Qual estratégia de tratamento?

\textbf{Q6.} $r = 0,82$ entre horas de sono e nota acadêmica. Você pode afirmar
que dormir mais \emph{causa} notas maiores? Justifique com dois argumentos.

\textbf{Q7.} Qual a diferença entre barras agrupadas e barras empilhadas 100\%?
Dê um exemplo de quando usar cada uma.

\textbf{Q8.} Um analista calcula a média de tempo de atendimento incluindo um
outlier de 8 horas (os demais variam entre 8 e 15 minutos). Como isso afeta a
média? O que você faria antes de reportar?

\textbf{Q9.} Você quer comparar o NPS (\textit{Net Promoter Score}, escala 1–10)
de dois grupos de clientes: Premium e Básico. Que medidas de posição e dispersão
você calcularia? Que visual usaria?

\textbf{Q10.} Um dataset de 200 pedidos tem 3 valores em branco na coluna
``data de entrega''. Os pedidos faltantes são de clientes aleatórios, sem relação
com região ou plano. Que tipo de missing? O que você faz?
\end{BoardBox}

\begin{SolvedBox}
\textbf{Gabaritos do Simulado (professor):}

\textbf{Q1.} Assimetria positiva intensa (cauda de altos salários puxa a média para cima).
Reportar mediana (R\$4.500 — mais representativa da maioria).

\textbf{Q2.} Mesma média, performances completamente diferentes. Turma A: homogênea,
processo consistente. Turma B: enorme variação — alguns alunos excelentes, outros muito abaixo.

\textbf{Q3.} Ordenados: $5,7,8,10,45$. $Q_1=6$; $Q_3=27{,}5$; $\text{IQR}=21{,}5$.
$\text{LS} = 27{,}5 + 32{,}25 = 59{,}75$. 45 $<$ 59,75 $\Rightarrow$ 45 \textbf{não é outlier} pela regra de Tukey neste conjunto.

\textbf{Q4.} Distribuição bimodal: dois subgrupos distintos no dataset (ex.: atendimentos
simples vs. complexos). Analisar cada subgrupo separadamente.

\textbf{Q5.} MNAR (ausência relacionada ao valor da variável em si — renda baixa omite renda).
Estratégia: criar flag \texttt{renda\_ausente} + imputar pelo percentil baixo do grupo com
padrão similar (ex.: escolaridade baixa).

\textbf{Q6.} Não. (1) Causalidade reversa possível: pessoas saudáveis dormem mais
\emph{e} têm melhor desempenho acadêmico — saúde pode causar os dois. (2) Variável
confundidora: gestão de tempo—quem gerencia bem o tempo dorme em horas regulares
\emph{e} estuda melhor.

\textbf{Q7.} Agrupadas: comparar valores absolutos de subcategorias lado a lado.
Ex.: receita de Produto A vs B por trimestre. Empilhadas 100\%: mostrar composição
proporcional. Ex.: share de mercado de cada canal por região.

\textbf{Q8.} A média é distorcida para cima — inutilizável como SLA. Antes de reportar:
identificar se 8h é erro de dado ou ocorrência real; se real, reportar mediana + mencionar
o outlier separadamente com contexto.

\textbf{Q9.} Posição: mediana por grupo (NPS pode ter outliers e assimetria).
Dispersão: IQR. Visual: boxplot lado a lado (Premium × Básico) — compara
distribuição completa, não apenas a média.

\textbf{Q10.} MCAR (ausência sem relação com nenhuma variável). Com apenas 1,5\% de
missing (3/200), pode remover as 3 linhas sem impacto material na análise.
\end{SolvedBox}

% ---------------------------------------------------------------------------
\section{Orientações finais — Prova I}
% ---------------------------------------------------------------------------

\begin{NoteBox}
\textbf{Para a Prova I (26/03):}
\begin{itemize}
  \item Fórmulas fornecidas: média, mediana, variância, DP, IQR, limites de Tukey, Pearson.
  \item Não é necessário memorizar fórmulas — é necessário saber interpretá-las e escolher
  a medida certa para cada contexto.
  \item Questões discursivas exigem: número + o que significa + limitação/risco + recomendação.
  \item Python: pode aparecer uma saída de \texttt{df.describe()} ou \texttt{df.corr()} para
  interpretação — não pedem para escrever código.
  \item Calculadora simples é permitida.
\end{itemize}
\end{NoteBox}

% ---------------------------------------------------------------------------
\section{Materiais Preparados}
% ---------------------------------------------------------------------------
\begin{itemize}
  \item Slides: resumo visual de toda a Unidade I (1 slide por tópico)
  \item Tabela 1 impressa para a prática
  \item Simulado impresso (1 folha por aluno)
  \item Gabarito do simulado para exibição após correção coletiva
\end{itemize}

% ---------------------------------------------------------------------------
\section{Próximo passo}
% ---------------------------------------------------------------------------

Após a Prova I, abrimos a \textbf{Unidade II: Análise Explanatória e
Storytelling com Dados}. A pergunta muda de \emph{``o que os dados mostram?''}
para \emph{``como eu conto essa história para que alguém tome uma decisão?''}

Livro novo: \textit{Datastory} (Nancy Duarte). A \textbf{Semana 08} (20/04)
traz as Apresentações do Trabalho 4 + a estrutura da Big Idea e do Storyboard.
